% Actividad local 7 = foro 3.2. Documento autónomo: no modifica las actividades 6, 8, 9 y 10.
\documentclass[12pt,letterpaper,oneside]{article}
\usepackage[margin=2.35cm]{geometry}
\usepackage[utf8]{inputenc}
\usepackage[T1]{fontenc}
\usepackage[spanish,es-tabla]{babel}
\usepackage[scaled=.96]{helvet}
\renewcommand{\familydefault}{\sfdefault}
\usepackage{setspace}
\usepackage{graphicx}
\usepackage{xcolor}
\usepackage{enumitem}
\usepackage{hyperref}
\definecolor{itescaRedDark}{HTML}{8F1117}
\definecolor{itescaInk}{HTML}{181818}
\hypersetup{colorlinks=true,linkcolor=itescaRedDark,urlcolor=itescaRedDark,
  pdftitle={3.2 Foro Toma de Decisiones},pdfauthor={Martín Jonathan de la Cruz Muñoz}}
\setlength{\parindent}{0pt}
\setlength{\parskip}{0.65em}
\setlength{\emergencystretch}{2em}
\onehalfspacing

\begin{document}
\color{itescaInk}
\begin{center}
  \includegraphics[width=4.2cm]{ITESCA/assets-itesca/web/logo-itesca-contorno.png}\\[0.3cm]
  {\large Instituto Tecnológico Superior de Cajeme\\
  Maestría en Gestión Administrativa}\\[0.3cm]
  {\LARGE\bfseries\color{itescaRedDark} 3.2 Foro Toma de Decisiones}\\[0.2cm]
  {\large Fundamentos de Gestión Administrativa · GGFG02\\
  Unidad III · Ponderación: 5\% · Actividad local 7}
\end{center}

\textbf{Estudiante:} Martín Jonathan de la Cruz Muñoz.\\
\textbf{Fecha de preparación:} 12 de septiembre de 2026.\\
\textbf{Vencimiento indicado:} 31 de agosto de 2026, 23:59.

\textbf{Estado:} participación inicial y dos respuestas personalizadas preparadas
para revisión. No publicadas durante esta preparación. Moodle muestra un aviso
de plazo vencido; debe comprobarse la aceptación de participaciones extemporáneas.
Este documento no constituye un comprobante de entrega.

\section*{Consigna y alcance}
Identificar la fase del proceso más propensa a errores y una de las fases
consideradas más importantes, justificar la postura y ofrecer recomendaciones
para mejorar la toma de decisiones gerenciales. Responder a dos compañeros.
La consigna y las aportaciones de referencia se consultaron en el
\href{https://cursos3.e-itesca.edu.mx/mod/forum/view.php?id=7794}{foro 3.2 de Moodle}
el 12 de septiembre de 2026. Los apartados siguientes contienen los textos propuestos.

\section{Participación inicial}
\textbf{Asunto: Diagnosticar con evidencia y revisar resultados para decidir mejor}

Buenas tardes, compañeras, compañeros y docente:

\textbf{Fase que considero más propensa a errores: la inteligencia o identificación
y análisis del problema.} Mi elección se debe a que, en esta etapa, se decide qué
situación necesita atención y cómo se interpreta la información. Confundir un
síntoma con su causa puede conducir a una solución bien ejecutada, pero dirigida
al problema equivocado. También existe el riesgo de buscar solamente datos que
confirmen una explicación previa o de atribuir una falla a una persona sin revisar
el proceso. Canós Darós y colaboradores (s.\,f., sección 4) señalan precisamente
el riesgo de responder adecuadamente a un problema mal formulado. No afirmo que
esta fase concentre siempre más errores; la considero especialmente vulnerable
por sus efectos sobre las etapas posteriores.

Por ejemplo, en un caso hipotético, una empresa que recibe quejas por entregas
tardías podría contratar más repartidores suponiendo que falta personal. Sin
embargo, el retraso podría originarse en registros de inventario incorrectos o
en autorizaciones internas. Antes de decidir, convendría revisar tiempos por
etapa, pedidos afectados y observaciones de quienes realizan el trabajo. Así se
distingue el problema observado de las posibles causas, que todavía deben verificarse.

\textbf{Una de las fases más importantes: la revisión o evaluación de resultados.}
Aunque un diagnóstico adecuado es indispensable, destaco la revisión porque permite
comprobar si la decisión resolvió el problema y detectar efectos no previstos. En
el modelo de Canós Darós y colaboradores (s.\,f.), esta fase evalúa lo realizado y
permite reiniciar el proceso cuando los resultados no son adecuados. Ejecutar una
decisión no equivale a demostrar su eficacia.

En el ejemplo anterior, propondría comparar el porcentaje de entregas a tiempo
antes y después de la intervención, sin perder de vista el costo por pedido y las
horas extra. Como meta ilustrativa, podría buscarse pasar de 80\% a 95\% de
puntualidad en cuatro semanas. Estas cifras son una propuesta, no datos de una
empresa real. También habría que revisar cambios en la demanda para no atribuir
toda mejora a la decisión tomada.

\textbf{Recomendaciones:} formular el problema sin anticipar la solución;
contrastar registros con la experiencia del personal; comparar alternativas por
costo, viabilidad, riesgos e impacto en las personas; asignar un responsable y un
plazo; y definir desde el inicio cómo se evaluarán los resultados. López, Guamán
y Castro (2020) destacan la relevancia de la información confiable y oportuna para
reducir la incertidumbre gerencial. Esto no significa esperar información perfecta,
sino explicitar lo que se sabe, lo que se supone y lo que falta comprobar.

En conclusión, considero que decidir mejor exige cuidar tanto el diagnóstico como
la revisión: el primero orienta la acción y la segunda permite corregirla y aprender
de sus resultados. ¿Qué indicador elegirían para comprobar que una decisión
solucionó la causa del problema y no solamente su manifestación?

\subsection*{Referencias de la participación}
  Canós Darós, L., Pons Morera, C., Valero Herrera, M., y Maheut, J. P. (s.\,f.).
  \textit{Toma de decisiones en la empresa: proceso y clasificación}.
  Universidad Politécnica de Valencia. Sección 4, «Etapas en el proceso de toma
  de decisiones». Lectura disponible en los materiales locales de la unidad III.

  López, D. A., Guamán, M. D., y Castro, J. C. (2020).
  La toma de decisiones y la eficacia organizativa en las PyMEs comerciales de la
  ciudad de Ambato (Ecuador). \textit{Revista Espacios, 41}(22), artículo 27.
  Secciones 1 y 1.1. Lectura disponible en los materiales locales de la unidad III.

\section{Respuesta a Mariel Nava Ortiz}
\textbf{Destino:}
\href{https://cursos3.e-itesca.edu.mx/mod/forum/discuss.php?d=615}{TOMA DE DECISIONES},
publicación inicial del 31 de agosto de 2026, 20:06.

\textit{Contexto verificado, no forma parte del mensaje:} Mariel considera central
la definición del problema y ubica el mayor riesgo de error en la recolección de
información por su duración, exigencias de análisis y costos.

\subsection*{Texto de la respuesta}
Hola, Mariel:

Me parece relevante que relaciones la recolección de información con el tiempo y
el costo de decidir. Coincido en que recopilar datos sin un objetivo claro puede
retrasar la solución. Añadiría que el riesgo no depende únicamente de la cantidad
de información, sino de su pertinencia y de cómo se interpreta. Incluso una base
amplia puede inducir a errores si solo incluye datos que respaldan lo que ya se pensaba.

Para complementar tu propuesta, establecería desde el inicio qué pregunta debe
responder cada dato, quién puede verificarlo y cuándo se necesita. En un problema
de entregas tardías, por ejemplo, sería más útil localizar en qué etapa se acumula
el retraso que reunir todos los reportes de la empresa. Esto permitiría delimitar
la búsqueda sin renunciar a contrastar explicaciones alternativas.

También coincido contigo en la importancia de definir el problema: un diagnóstico
explícito orienta qué información resulta útil. ¿Qué criterio utilizarías para
decidir que ya se cuenta con información suficiente para actuar, sin prolongar
el análisis ni asumir riesgos innecesarios?

\section{Respuesta a Laura América González Valenzuela}
\textbf{Destino:}
\href{https://cursos3.e-itesca.edu.mx/mod/forum/discuss.php?d=619}{Toma de decisiones},
publicación inicial del 31 de agosto de 2026, 22:58.

\textit{Contexto verificado, no forma parte del mensaje:} Laura identifica errores
en el diagnóstico, distingue causas de efectos y destaca la evaluación de alternativas
según recursos, riesgos, beneficios y consecuencias. Recomienda escuchar al personal
y revisar resultados.

\subsection*{Texto de la respuesta}
Hola, Laura:

Coincido con tu distinción entre causas y efectos y con la necesidad de no elegir
automáticamente la primera alternativa. Tu propuesta de valorar recursos, riesgos
y consecuencias podría concretarse mediante una matriz de comparación cuyos criterios
se acuerden antes de calificar las opciones. Esto ayudaría a evitar que los criterios
se acomoden a una solución preferida de antemano.

Añadiría que una alternativa con el menor costo inmediato no necesariamente es la
más conveniente. Por ejemplo, ante entregas tardías, recurrir permanentemente a horas
extra podría parecer más rápido que corregir el inventario, pero también generar
fatiga y nuevos errores. Por eso incluiría el impacto en el personal y la sostenibilidad
de la solución, además del costo y el plazo.

Tu recomendación de evaluar después de implementar complementa mi postura sobre la
revisión: la comparación anticipada necesita contrastarse con resultados reales.
¿Cómo priorizarías dos alternativas si una ofrece resultados rápidos y la otra
reduce mejor el riesgo de que el problema se repita?

\section*{Control de preparación y publicación}
\begin{itemize}[leftmargin=*]
  \item \textbf{Preparado:} postura justificada sobre ambas fases, recomendaciones,
  ejemplo hipotético, referencias consultadas y dos respuestas a mensajes reales.
  \item \textbf{Por realizar:} revisión y aprobación del estudiante, comprobación
  de aceptación de participación extemporánea y publicación de los tres mensajes.
  \item \textbf{Evidencia de entrega:} conservar los enlaces y fechas reales de
  las tres publicaciones una vez enviadas. No se han generado durante esta preparación.
\end{itemize}
El PDF es un respaldo; no sustituye las intervenciones en Moodle. No se ha pulsado
«Marcar como hecha» ni se afirma que haya ocurrido un intercambio posterior con
las compañeras. La versión editable para copiar cada mensaje se conserva en el
documento Markdown de participación de la Actividad 7.
\end{document}
